% ============================================================================
% CAPÍTULO 8: NEUTRINOS
% ============================================================================
\chapter{Neutrinos}
\label{cap:neutrinos}

\begin{center}
\textit{``Las partículas fantasma''}
\end{center}

\vspace{0.5cm}

\begin{tcolorbox}[colback=white, colframe=kleinblue, boxrule=2pt]
\centering
\textit{``Puedo calcular el movimiento de los cuerpos celestes,\\
pero no la locura de las masas.''}\\
--- Isaac Newton\\[0.3cm]
\textit{``Yo puedo calcular ambas cosas: $m_\nu = m_e / 2(7\pi)^5$.''}\\
--- Teoría Klein, 2026
\end{tcolorbox}


% ============================================================================
\section{Las Partículas Más Esquivas}
% ============================================================================

Los neutrinos son las partículas más abundantes del universo después de los fotones.
Pero casi no interactúan con la materia:

\begin{itemize}
    \item $10^{14}$ neutrinos solares atraviesan tu cuerpo cada segundo
    \item Casi todos pasan sin tocar un solo átomo
    \item Se necesitan detectores de kilómetros de tamaño para capturar algunos
\end{itemize}

\textbf{EL MISTERIO DE LA MASA:}

Durante décadas se pensó que los neutrinos no tenían masa.
Hoy sabemos que sí tienen, pero es MUY pequeña:
\begin{equation}
m_\nu < 0.1 \text{ eV} \quad \text{vs} \quad m_e = 511,000 \text{ eV}
\end{equation}

¿Por qué son tan livianos?


% ============================================================================
\section{La Masa del Neutrino: $m_e/m_{\nu_3} = 2 \times (7\pi)^5$}
% ============================================================================

\begin{nivelmatematico}
\textbf{Masa del Neutrino desde Klein}

\begin{equation}
\boxed{\frac{m_e}{m_{\nu_3}} = 2 \times (7\pi)^5}
\end{equation}

donde $\nu_3$ es el neutrino más pesado.

\textbf{Desglose:}
\begin{itemize}
    \item $(7\pi)^5 \approx 5.17 \times 10^6$
    \item $2 \times (7\pi)^5 \approx 1.03 \times 10^7$
    \item $m_{\nu_3} = m_e / (2 \times (7\pi)^5) \approx 50$ meV
\end{itemize}

\textbf{VERIFICACIÓN:}
\begin{align}
\text{Predicción:} & \quad m_{\nu_3} \approx 50 \text{ meV}\\
\text{Observado:} & \quad m_{\nu_3} \approx 50-60 \text{ meV (cosmología)}\\
\text{Error:} & \quad \sim 1.2\%
\end{align}

\textbf{INTERPRETACIÓN:}

Los neutrinos son livianos porque su masa está suprimida por $(7\pi)^5$.
El exponente 5 = dimensiones de Kaluza-Klein.
El factor 2 corresponde a dos quiralidades (izquierda/derecha).
\end{nivelmatematico}


% ============================================================================
\section{El Ángulo de Mezcla: $\theta_{13} = 1/7$ rad}
% ============================================================================

Los neutrinos ``oscilan'' entre tres sabores ($\nu_e$, $\nu_\mu$, $\nu_\tau$).
Los ángulos de mezcla determinan estas oscilaciones.

\begin{nivelmatematico}
\textbf{FÓRMULA KLEIN:}

\begin{equation}
\boxed{\theta_{13} = \frac{1}{7} \text{ rad} \approx 8.19°}
\end{equation}

\textbf{VERIFICACIÓN:}
\begin{align}
\text{Predicción:} & \quad \theta_{13} = 8.19°\\
\text{Observado:} & \quad \theta_{13} = 8.54° \pm 0.15°\\
\text{Error:} & \quad 4.0\%
\end{align}

\textbf{INTERPRETACIÓN:}

El factor $1/7$ aparece porque este ángulo conecta la primera y tercera generación,
``saltando'' una capa de la estructura Klein de 7 capas.
\end{nivelmatematico}


% ============================================================================
\section{Predicciones para Experimentos Futuros}
% ============================================================================

\begin{nivelconceptual}
\textbf{Experimentos que pueden verificar Klein:}

\begin{enumerate}
    \item \textbf{KATRIN} (Alemania): Mide $m_\nu$ directamente del decaimiento beta.
    Predicción Klein: $m_{\nu_3} \approx 50$ meV.

    \item \textbf{Project 8} (EEUU): Nueva técnica de medición ultra-precisa.
    Puede alcanzar sensibilidad de 40 meV.

    \item \textbf{DUNE} (EEUU): Medirá ángulos de mezcla con alta precisión.
    Verificará $\theta_{13} = 1/7$ rad.

    \item \textbf{JUNO} (China): Determinará la jerarquía de masas de neutrinos.
\end{enumerate}
\end{nivelconceptual}


% ============================================================================
\section{Resumen del Capítulo}
% ============================================================================

\begin{tcolorbox}[colback=kleingray, colframe=kleinblue, title=\textbf{Resumen del Capítulo 8}]

\textbf{PREDICCIONES NEUTRINOS KLEIN:}

\begin{enumerate}
    \item $m_e/m_{\nu_3} = 2 \times (7\pi)^5$ \hfill [1.2\% error] --- Masa del neutrino
    \item $\theta_{13} = 1/7$ rad \hfill [4.0\% error] --- Ángulo de mezcla
\end{enumerate}

\vspace{0.3cm}
\textbf{PREDICCIÓN NUMÉRICA:}

\begin{equation}
m_{\nu_3} \approx 50 \text{ meV}
\end{equation}

Esta predicción será verificable en la próxima década con KATRIN y Project 8.

\end{tcolorbox}


% ============================================================================
\section*{Código de Verificación}
% ============================================================================

\begin{lstlisting}[caption={Verificación numérica - Capítulo 8}]
import numpy as np
from numpy import pi

siete_pi = 7 * pi

# Masa del electron en eV
m_e_eV = 511000  # eV

# Masa del neutrino
ratio_klein = 2 * siete_pi**5
m_nu_klein = m_e_eV / ratio_klein  # en eV
print(f"MASA DEL NEUTRINO:")
print(f"  m_e/m_nu Klein = {ratio_klein:.2e}")
print(f"  m_nu Klein = {m_nu_klein*1000:.1f} meV")
print(f"  m_nu obs   ~ 50-60 meV")

# Angulo de mezcla theta_13
theta_13_klein = 1/7  # rad
theta_13_deg_klein = np.degrees(theta_13_klein)
theta_13_obs = 8.54  # grados
print(f"\nANGULO DE MEZCLA theta_13:")
print(f"  theta_13 Klein = {theta_13_deg_klein:.2f} grados")
print(f"  theta_13 obs   = {theta_13_obs:.2f} grados")
print(f"  Error = {abs(theta_13_deg_klein - theta_13_obs)/theta_13_obs*100:.1f}%")
\end{lstlisting}
